\subsection{The Spacetime Partition (Weak Mixing Angle \texorpdfstring{$\sin^2 \theta_W$}{sin2thetaW})}
\label{subsec:spacetime_partition}

\CatchFileBetweenTags{\AlphaInvVal}{constants.tex}{AlphaInvVal}
\CatchFileBetweenTags{\WeakAngleStepOneVal}{constants.tex}{WeakAngleStepOneVal}
\CatchFileBetweenTags{\WeakAngleVal}{constants.tex}{WeakAngleVal}
\CatchFileBetweenTags{\WBosonMassVal}{constants.tex}{WBosonMassVal}
\CatchFileBetweenTags{\CompTOL}{constants.tex}{CompTOL}
\CatchFileBetweenTags{\WeakAngleTCheckVal}{constants.tex}{WeakAngleTCheckVal}
\CatchFileBetweenTags{\WeakAngleExperimentalValue}{constants.tex}{WeakAngleExperimentalValue}
\CatchFileBetweenTags{\WeakAngleGlobalExperimentalValue}{constants.tex}{WeakAngleGlobalExperimentalValue}
\CatchFileBetweenTags{\WeakAngleGlobalSigma}{constants.tex}{WeakAngleGlobalSigma}

\CatchFileBetweenTags{\NcVal}{constants.tex}{NcVal}
\CatchFileBetweenTags{\NcEq}{constants.tex}{NcEq}
\CatchFileBetweenTags{\WeinbergGUTVal}{constants.tex}{WeinbergGUTVal}
\CatchFileBetweenTags{\WeinbergGUTEq}{constants.tex}{WeinbergGUTEq}

\subsubsection{The Architectural Requirement}

Under \textbf{Rule 1: Capacity Partitioning}, a persistent system must geometrically divide its finite processing bandwidth between maintaining its structural state and updating that state. This manifests as the \textbf{Spacetime Partition Ratio}: the exact fraction of the vacuum's total informational bandwidth allocated strictly to temporal evolution (the fundamental causal update rate, $\Delta$) versus spatial and internal structural maintenance. In the Standard Model, this exact architectural allocation is observed as the Weak Mixing Angle ($\sin^2 \theta_W$). 

\subsubsection{The Total System Capacity (\texorpdfstring{$N_{sys}$}{Nsys})}
To derive the partition, we must quantify the total active bandwidth of the lattice. The \textbf{Total System Capacity} ($N_{sys}$) is the sum of the substrate's three independent geometric primitives (Space, Matter, and Force) that must share the multiplexed manifold:

\begin{enumerate}
    \item \textbf{Spacetime Structure ($R_M$):} The fundamental lattice resonance ($\Delta=43$) must be embedded across all $D=4$ spacetime dimensions to define the propagation metric. This is the Manifold Resolution ($R_M$) defined previously:
    \begin{equation}
        R_M = D \times \Delta = 4 \times 43 = 172
    \end{equation}
    
    \item \textbf{Matter Content ($\nu \cdot \eta$):} The chiral spinor manifold requires $\nu=16$ degrees of freedom to encode fermion quantum numbers. This acts as a \textbf{bulk capacity} that must be distributed across the discrete lattice volume ($R_M = 172$ nodes) and subject to the Coordinate Overhead ($\eta$). (Note $R_M$ defines the geometric volume itself and $C_{force}$ is an algebraic rank (a boundary/topological invariant); therefore, only $C_{matter}$, which must be distributed across the bulk manifold, is subject to $\eta$). 
    \begin{equation}
        C_{matter} = \nu \cdot \eta \approx 15.9069
    \end{equation}
    
    \item \textbf{Interaction Rules ($\sigma$):} The interaction group possesses $\sigma=5$ independent generators. (\textit{Note: While the Standard Model features 12 total gauge bosons, the underlying geometric \textbf{rank}, the number of mutually commuting generators required to span the interaction space is strictly $\sigma=5$}).
    \begin{equation}
        C_{force} = \sigma = 5
    \end{equation}
\end{enumerate}


These three orthogonal sectors add linearly to define the Total System Capacity:
\begin{equation}
\label{eq:system_capacity}
\begin{split}
N_{sys} &= R_M + C_{matter} + C_{force} \\
&= 172 + 15.9069 + 5 \\
&\approx \mathbf{\WeakAngleStepOneVal}
\end{split}
\end{equation}

\subsubsection{The Partition Formula}
The Weak Mixing Angle is defined as the strict ratio of the \textbf{Temporal Resonance} (the bare causal frequency $\Delta$) to the \textbf{Total System Capacity} ($N_{sys}$). The numerator is the bare temporal resonance $\Delta = 43$ (the number of discrete causal slots per fundamental cycle), not the temporal dimensionality ($D_t = 1$), because the partition allocates \textit{capacity} (state slots), not \textit{dimensions}.
\begin{equation}
\sin^2 \theta_W = \frac{\text{Temporal Resonance}}{\text{Total Capacity}} = \frac{\Delta}{N_{sys}}
\end{equation}
Substituting the invariants:
\begin{equation}
\sin^2 \theta_W = \frac{43}{\WeakAngleStepOneVal} \approx \mathbf{\WeakAngleVal}
\end{equation}

\vspace{0.5em}
\begin{tabular}{ll}
\textbf{Prediction:} & \num{\WeakAngleVal} \\
\textbf{Experiment (Direct Mass):} & \WeakAngleExperimentalValue \\
\textbf{Experiment (Global Fit):} & \WeakAngleGlobalExperimentalValue \\
\end{tabular}
\vspace{0.5em}

The Weak Mixing Angle is exactly the bandwidth fraction allocated to temporal updates ($\Delta$). Because the lattice strictly partitions physical on-shell capacity, the geometry inherently aligns with the Direct Mass measurement.

The deviation from the Global Fit (\WeakAngleGlobalSigma $\sigma$) confirms that the geometry does not require the continuous, infinite loop corrections assumed by the Standard Model fit, directly predicting the empirical tension observed in the W-Boson mass.

\subsubsection{Physical Interpretation: Electroweak Structure}
In the Standard Model, the weak mixing angle relates the gauge couplings of the $U(1)_Y$ hypercharge ($g'$) and the $SU(2)_L$ weak isospin ($g$) via the relation: $\sin^2\theta_W = g'^2 / (g^2 + g'^2)$.

The $E_8$-Persistence theory establishes a structural isomorphism between these abstract couplings and the lattice geometry:
\begin{equation}
\frac{g'^2}{g^2 + g'^2} \longleftrightarrow \frac{\Delta}{R_M + \nu\eta + \sigma}
\end{equation}
This geometric partition ($\approx 22\%$ Temporal vs. $78\%$ Spatial) dictates the entire phenomenology of the electroweak gauge bosons:
\begin{itemize}
    \item \textbf{The Photon ($\gamma$, $U(1)$ Temporal):} The hypercharge sector ($g'$) couples strictly to the \textbf{Fundamental Resonance} ($\Delta$), representing the causal update rate of the lattice. Because temporal evolution cannot be obstructed without violating causality, the photon must remain massless and propagate at $c$.
    \item \textbf{The W/Z Bosons ($W^{\pm}, Z^0$, $SU(2)$ Spatial):} The isospin sector ($g$) couples to the \textbf{Manifold Embedding} ($N_{sys} - \Delta$), representing the dense spatial infrastructure. To propagate through this structural density, these gauge fields must deform the lattice geometry, incurring a severe geometric impedance cost that manifests as mass.
\end{itemize}

\subsubsection{Validation: Resolution of the W Boson Mass Tension}
While the absolute electroweak mass scale is derived in System 5 (The Resonant Interface), we can immediately validate the geometric partition ratio by testing the relative scaling between the $W$ and $Z$ bosons. The most stringent test of this geometric partition is the mass of the W Boson. The Standard Model prediction for $M_W$ is currently in tension with precise measurements.

\begin{itemize}
    \item \textbf{Standard Model:} $80.357 \pm 0.006$ GeV.
    \item \textbf{CDF II (2022):} $80.4335 \pm 0.0094$ GeV ($7\sigma$ tension).
    \item \textbf{ATLAS (2023):} $80.360 \pm 0.016$ GeV (SM-consistent).
\end{itemize}

Using the derived angle $\sin^2 \theta_W \approx 43/193$ and the experimental Z-boson mass ($M_Z = 91.1876$ GeV), we calculate the W mass geometrically:
\begin{equation}
\begin{split}
    M_W &= M_Z \sqrt{1 - \sin^2 \theta_W} \\
    &= 91.1876 \times \sqrt{1 - \frac{43}{\WeakAngleStepOneVal}} \\
    &\approx \boxed{\mathbf{\WBosonMassVal}}
\end{split}
\end{equation}

\textbf{Result:} The geometric prediction ($M_W = \WBosonMassVal$) lies between the Standard Model ($80.357$ GeV) and CDF II ($80.434$ GeV), validating the \textbf{direction} of the CDF anomaly (favoring a heavier W-boson) while remaining consistent with ATLAS ($80.360$ GeV) within $1.6\sigma$. The lattice geometry predicts a heavier W than the SM global fit through a rigid structural partition ($\approx 43/193$) rather than parameter adjustment, suggesting the SM value may be systematically underestimated.

\textit{Note on Radiative Corrections:} In standard QFT, calculating $M_W$ directly from $M_Z$ and $\sin^2\theta_W$ requires complex perturbative loop corrections ($\Delta r$) because the bare Lagrangian parameters are unphysical. In Informational Energetics, because $\sin^2\theta_W \approx 43/193$ represents the literal, physical partition of the lattice bandwidth, it acts as the fully dressed, non-perturbative geometric ratio. Therefore, the exact geometric relation holds for the on-shell masses without requiring perturbative loop approximations. If precision measurements converge on a value of $M_W$ that differs from $M_Z\sqrt{1 - 43/N_{sys}}$ by more than the experimental uncertainty, it would falsify the identification of $\sin^2\theta_W$ as the exact physical partition ratio.

\subsubsection{Implications for Grand Unification: The Symmetric Precursor}

Standard Grand Unified Theories (such as $SU(5)$) predict that at the unification scale, before symmetry breaking, the Weak Mixing Angle converges to $\sin^2 \theta_W = 3/8 = 0.375$. 

The $E_8$-Persistence framework recovers this exact theoretical value naturally, identifying it not as a dynamic high-energy convergence, but as the \textbf{Unprojected Symmetric State} of the lattice. 

Before the lattice undergoes Chiral Truncation to establish the Arrow of Time (which isolates the anti-chiral sector and projects the geometry down to $D=4$), both $D_4$ orthogonal sublattices of the $E_8 \to D_4 \oplus D_4$ decomposition are fully active, yielding a total manifold rank of $2D = 8$. In this symmetric phase, the bandwidth partition is simply the ratio of the Color Sector ($N_c = \NcEq = \NcVal$) to the Total Lattice Rank:
\begin{equation}
    \sin^2 \theta_W(M_{GUT}) = \WeinbergGUTEq = \mathbf{\WeinbergGUTVal}
\end{equation}

This demonstrates a profound architectural relationship. The low-energy Weak Mixing Angle derived previously ($\sin^2 \theta_W = 43/193 \approx 0.223$) is the exact, physical result of projecting this symmetric $3/8$ ratio onto the constrained, time-asymmetric 4D causal manifold. 

Standard GUTs fail to reach the correct low-energy values (predicting $\approx 0.21$ at the Z-pole rather than the observed $0.223$) because they assume the continuous $3/8$ symmetry dynamically breaks and runs via continuous renormalization. The $E_8$-Persistence theory resolves this: the apparent ``non-unification'' of gauge couplings is evidence that forces arise from distinct geometric features ($\Delta$ vs. $D\Delta$ vs. $\sigma$) rather than a single broken symmetry group, making $\sin^2\theta_W$ a \textbf{fixed geometric ratio} dictated by projection, not a running coupling.


\subsubsection{Recursive Validation: The Cost of Time}
The structural consistency of the theory is confirmed by linking this partition back to the Fine-Structure Constant. Recall the \textbf{Toll} component of the $\alpha^{-1}$ geometric impedance ($T_{geo}$), derived from the 3-point interaction vertex on the substrate in the companion paper\cite{meyer_informational_2026}:
\begin{equation}
\begin{split}
    T_{geo} &= \frac{1}{(2\nu)^3} \cdot \frac{\chi}{\sigma} \cdot \left(1 - \frac{\sigma}{D\Delta} \right) \approx \CompTOL
\end{split}
\end{equation}

We now observe that this bare geometric transition cost ($T_{geo}$) structurally mirrors the fully dressed, second-order electromagnetic cross-section ($\alpha^2$, as expected for a 3-point vertex) modulated by the Weak Mixing Angle ($\sin^2 \theta_W$) we just derived:
\begin{equation}
\begin{split}
T_{check} &= \alpha^2 \sin^2 \theta_W \\
&= \left(\frac{1}{\AlphaInvVal}\right)^2 \cdot \left(\frac{43}{\WeakAngleStepOneVal}\right) \\
&\approx (5.325 \times 10^{-5}) \cdot (0.2229053) \\
&\approx \WeakAngleTCheckVal
\end{split}
\end{equation}

\textbf{Conclusion:} The match to within $\sim 0.15\%$ confirms internal consistency across the macroscopic and microscopic derivations of the framework. The Toll $T_{geo}$ represents the \textit{bare} isolated geometric cost of a temporal vertex, while $T_{check}$ evaluates the \textit{fully dressed} macroscopic interaction utilizing the complete vacuum impedance ($\alpha$). That these two entirely independent geometric paths, one derived from 3-point vertex geometry, the other from macroscopic bandwidth partitioning, converge to nearly identical values confirms that Weak Interaction (time evolution) is strictly the temporal bandwidth allocation on the electromagnetic manifold.